% ============================================================
%  Harald Høffding, "Religion og Videnskab" — "Religion and Science"
%  Religionshistoriske Smaaskrifter I
%  Copenhagen: Gyldendal — Nordisk Forlag, 1910
%  TRANSLATION (English)
%  Mirrors transcription.tex 1:1. Source of truth = transcription.tex.
%  Progress: printed pp. 1–5 (PDF 10–14) translated; continues mid-work.
% ============================================================
\documentclass[12pt,a4paper]{article}
\usepackage[utf8]{inputenc}
\usepackage[T1]{fontenc}
\usepackage{libertinus}
\usepackage{libertinust1math}
\usepackage{textalpha}
\usepackage[english]{babel}
\usepackage[protrusion=true,expansion=true]{microtype}
\usepackage{setspace}
\usepackage[margin=1.2in]{geometry}
\usepackage{fancyhdr}
\usepackage{enumitem}
\usepackage{hyperref}

\hypersetup{
  pdftitle={Religion and Science},
  pdfauthor={Harald Høffding},
  hidelinks
}

\pagestyle{fancy}
\fancyhf{}
\rhead{Høffding, \emph{Religion and Science} (1910)}
\cfoot{\thepage}

\setstretch{1.3}

\title{\textbf{Religion and Science}\\[0.4em]
\large \emph{Religionshistoriske Smaaskrifter} I}
\author{Harald Høffding\\[4pt]
  \small Translated from the Danish}
\date{Copenhagen: Gyldendal --- Nordisk Forlag, 1910}

\begin{document}
\maketitle
\thispagestyle{fancy}

\bigskip

\noindent Within every culture we know, a critical condition sets in once
the division of labour comes into force. In the material sphere it then
becomes no longer possible for the individual worker to produce a work in
which he can have the feeling of having worked with all his powers. What
the individual can do now becomes, most often, only the production of a
single part of a whole, a part that requires merely the application of a
single faculty. The development of his power and his talent thereby becomes
one-sided. Like his work, his very personality readily becomes but a
fragment. In the spiritual sphere the division of labour means that art,
science, religion, and morality, which in earlier times formed a unity, now
step apart and often step into opposition to one another as contending
powers. In earlier cultural periods religion stood as the great central
current of life, to which all spiritual forces and drives gave their
contribution, and which on the other hand yielded full satisfaction for the
artistic and the intellectual need, in so far as such need stirred at all.
But gradually the spiritual needs seek their satisfaction along different
paths. The investigative urge is no longer satisfied by the answers that
religious ideas can give to the manifold questions that life and experience
may raise. Special methods are developed, special forms for posing the
questions and for seeking the answers. The artistic impulse likewise goes
more and more its own ways, and an independent foundation is sought for
moral valuation. History shows that, in contrast to these new formations,
religion has a tendency to hold fast to the notions and modes of action
once established. It resists the spiritual division of labour as long as it
can, and it seeks as long as possible to close its eyes to the altered
conditions of spiritual life. From this follows a dividedness and a
splintering in spiritual life. There is lacking that gathering and harmony,
that concentration of all the forces of the mind in one common direction,
which under the earlier conditions was comparatively easy to attain. On
this rests the fact that there exists a religious problem, and that such a
problem will continue to exist until the spiritual life of individuals and
of the race as a whole has won new gathered and gathering forms for the
satisfaction of its innermost need---if indeed such new forms are possible
at all.

I shall here bring forward only a few features of the intellectual side of
the problem. The whole question is, for me, not exhausted by being viewed
from this side. But it is of great importance to become clear about what
the opposition that again and again asserts itself between religion and
science really rests upon. Often one finds---and this holds of both parties
that here stand over against one another---the ground of this opposition
put in the wrong place. There is much superficial and dogmatic harping on
science's real or supposed results on the one side, and there is on the
other side much short-sightedness with respect to the change that the life
of thought has in fact undergone, quite apart from the many conspicuous
results. In the collisions to which the opposition leads, the onlooker does
not miss honesty in the confessions made from the one side or the other.
For nowadays there is, after all, comparatively little danger in making
one's confession, in whatever direction it may go. By contrast, one has
ample occasion to miss that intellectual integrity which does not let the
mind rest before the great questions have been examined from all sides and
by all the means that stand at one's disposal. This is unfortunately not a
virtue that religious ethics has especially inculcated, and many of
religion's opponents show, by the character of their polemic, that they too
do not know the strict demands it makes. The want that an onlooker may feel
here is bound up with the fact that religious debate in our days has mostly
the character of agitation. On both sides it is a matter of winning
adherents, of dissemination, not of deepening. New thoughts, accordingly,
this modern debate has not brought either.

What I wish to bring here is likewise not new thoughts, but old thoughts
that have not been allowed to come to their full due.

\begin{center}---\end{center}

If we go back to times in which religion was not merely help and
consolation, but also gave answers to what thought could then ask
about---times, that is, in which the whole conception of the world had a
religious character---then there are two peculiarities that especially come
to the fore. The whole conception of the world has an intuitive character.
Existence stands as a whole whose several regions can be seen or at least
represented, without its being necessary to break with what the senses
immediately seemed to teach. Cosmic space, as it presented itself to
sensation and imagination, contained the possibility of localizing all the
great events with which religious proclamation dealt. And when
changes---especially momentous changes---entered into this world of
intuitability, it was not by the arduous path of research that \emph{the
explanation} was to be sought. Such changes were explained by the
intervention of personal beings in the course of events. It is
intuitability, and it is the principle of personal causes (what in the
widest sense may be called animism), that characterizes a specifically
religious conception of the world.

But both the concept of space and the concept of cause have, under the
influence of the independent urge to knowledge, undergone changes that are
here of decisive importance.

\medskip

1. We have a natural urge to refer everything that happens to definite
places in space, and religion in its classical periods does not deny this
urge. Religious representations have the character of intuition, not of
concept, and their great influence, their power to grip and to fill, rests
in no small part upon this. The content of the higher popular religions is
a series of historical events into which the fate of men and of human life
is woven---a great drama, and a drama played out upon a stage,
sense-intuition, which imagination can in any case very well encompass
without bursting. The whole of existence, so far as it is accessible at
this stage, is taken up into this stage, and its various regions furnish
places for the various scenes in the drama. The earth stands as centre, or
as the lowest place, and from the visible heaven, which is the world of the
deity, divine beings descend to intervene in the fate of men, perhaps
themselves to suffer and die, before they ascend again into the heavenly
regions, to which they thereupon open access for those who will follow
them. ``All men,'' says Aristotle, ``suppose that there are gods, and all,
both Hellenes and barbarians, assign to the deity the highest place.'' This
utterance applies to all popular religions, however different they may
otherwise be. Religious representations were woven in a natural way into
men's whole involuntarily formed conception, and there as yet existed no
science that necessitated a radical change of that conception. There was no
vagueness or indistinctness in the representation of the divine; this could
be localized like everything else; one could point toward the place from
which the divine forces welled forth, as well as to the place at which they
worked. The divine thoughts were, to be sure, raised as high above man's
thoughts as heaven above earth; but heaven was precisely not raised so high
above earth that intuition had to give up. Despite the difference in power
and dignity, a vivid and intimate relation to the gods was nevertheless
possible.
No popular religion has abandoned this picture of the world. But even before
science had seriously shaken the intuitability of the world, misgivings did arise
among thinking minds, and the security with which men had hitherto dwelt upon the
intuitable world-picture as the frame for the whole content of religion was
shaken. And this unrest was due in essential part to a deepening of the religious
consciousness itself. There stirs in all religion a twofold urge: on the one hand
an urge to raise the object of the religious representations as high above human
and earthly conditions and forms as possible, and on the other hand precisely an
urge to feel oneself one with the deity, or at least to be near it. It is an urge
to stand before an Exalted One, and an urge toward an intimate relation to the
Highest. Along both paths a break with the naïve intuitability of popular
religion could be brought about. If the deity is to be exalted above all things,
must it not then also be exalted above all spatial distinctions, above the
possibility of localization and limitation? And if it is to stand in the most
intimate relation to man, can it then be separated in space from the human being
with whom it is to be one? A distance in space is always an external relation,
however quickly it may be traversed. In the original religious documents
themselves both tendencies are to be traced; there makes itself known a purely
religious urge that had to lead out beyond intuitability. With this there united
itself also the influence of philosophical reflection, in the case of
Christianity especially the influence of Platonism. It was under the working of
these various motives that Augustine could say of his God that he was
``everywhere and nowhere,'' and that he had his place ``further within the soul
than its own inmost being,'' as well as ``higher up than the soul knew how to
climb.'' But beforehand there was first a hard---both bodily and
spiritual---struggle between the religious materialists (anthropomorphists), who
held to absolute intuitability, and the religious spiritualists, who denied it.

In theory it was spiritualism that won. But men could not and would not wholly
deny intuitability, which after all appears so clearly and unambiguously in the
original religious documents and is moreover of such great pedagogical
importance. In a narrative such as that of Jesus' ascension it is indeed
altogether necessary, if it is to be taken literally, to conceive heaven as a
definite place, up to which the journey went. The ascension is a passage from the
earthly part of space to the heavenly. The theologians of the Middle Ages then
took the position of teaching that the ascension meant a passage to another kind
of existence, a spiritual mode of being. To ascend to heaven means to lay aside
earthly and human limitation. But this symbolizing conception did not for them
exclude the literal one. Men held fast at once to intuitability and to symbolism.
A passage has really taken place from physical earth to physical heaven, but this
passage itself means at the same time something else. The literal and the
figurative meaning of ``the Higher'' in contrast to ``the Lower'' are both held
fast---though men are evidently inclined to lay the greater weight on the
figurative. It was assumed that the twofold meaning---the realistic and the
idealistic, the literal and the edifying---had been laid down in the original
accounts themselves. What for the original religious consciousness stood as an
undivided truth, intuitable and significant at once and at a single stroke, stood
for the medieval consciousness as a Both---And, at once separated and inseparable.

In Dante's great poem this standpoint prevails. The symbolic meaning is indeed
here too the one that lies nearest the author's heart. But the ideal truth has a
corresponding outward garb, whose reality is held fast. The unique thing in
Dante's \emph{Divina Commedia} rests, seen from this point of view, precisely on
this: that the outer framework stands fast, has a massiveness and an
intuitability that the religious poetry of later ages could not attain, and that
nevertheless the inwardness comes fully into its own through the symbolic
interpretation that shines through. The passage from the material to the ideal is
made---naïvely and grandly at once---at the very boundary of space: the outermost
heaven, which encloses all things, has itself its ``Where?'' in the deity, whose
power and love is the source of the revolution of the worlds.

But there came the time when men demanded a more prosaic answer to the question
``Where?'' Just as, within our experience, we determine every place by its
relation to another place, so we must also assume that everywhere in the world,
wherever we may think ourselves, the same holds. No absolute places or absolute
boundaries can then be thought. Every intuitable object has its ``Where'' in the
literal sense, determined by relation to other places. Therefore no world-centre
and no world-boundary can be represented. No place can be found out from which
the divine forces might stream forth or down. At every place the same law, or
rather the same method for the determination of place, must find application. It
was precisely this train of thought that led Giordano Bruno to draw his bold
consequences from Copernicus' hypothesis. Already in Lucretius' famous didactic
poem this train of thought was hinted at, but only far later did it become
possible to hold it fast and apply it. Over against the immeasurable world of
space that now opened itself, the ancient contrast between heaven and earth had
to stand as vanishing. That space is in fact infinite cannot, in a certain sense,
be proved. The essential thing is the new method, the new way of thinking, when
the talk is of place and space. It is in a way a new organ that develops, and
that thought brings with it wherever in the world of space it may betake itself.
Over against this new way of thinking the old world-picture with its absolute
places was untenable. The passion with which it was so long held fast, and with
which men rushed upon the champions of the new world-picture, shows too that men
had an inkling of what was at stake. For what was at stake was nothing less than
the surrender of the literal meaning of the religious representations. All
popular religions had laid the old, absolute, and limited world-picture at their
foundation. It furnished the ground and the point of support for the religious
imagination and for the feeling of at-homeness with which the human spirit could
regard a world in which its highest ideals could have their ``natural places''
just as well as the meanest material elements had theirs.

The medieval Both---And was now no longer sufficient. The events of which the
biblical writings tell so reliably and intuitably must now, at decisive points,
be understood figuratively. One can no longer help oneself by assuming both a
real and an ideal meaning. The narrative of the ascension is here the typical
example, because in it it is clearly expressed that the whole universe is
affected by what happened in Palestine. A heavenly power had descended to earth
and now left it again to return to its natural place.

Only very seldom does the religious consciousness seem to enter vividly into the
questions that hereby arise, and that are of importance both for the historical
understanding of the biblical narratives and for the characterization of the
religious world-view. If the narrative of the ascension cannot be literally true,
then it stands as a legend that may have its beauty and its ideal significance
when we place ourselves back into the thought of the time, but that can no longer
come forward with dogmatic weight. Men have thereby, on an essential point,
acknowledged the legitimacy of the historical-critical method. For this method
rests on the same presupposition on which Bruno's extension of the method of
determining place from the earth to heaven rested, namely that in the
investigation of the distant we can do nothing but apply the procedure that we
have tried upon the near. We must meet the old accounts with the world-conception
that we now have, even though it be ever so different from that which underlies
those accounts. But it is not only for the historical method that the
consequences become important. The whole significance that can be attached to
religious representations must thereby become a different one. No wonder that men
prefer to let the matter rest! And yet even orthodox theologians here
involuntarily become symbolists and rationalists, as soon as they seriously think
their way into the subject. Martensen, for example, maintains firmly that the
description of the ascension is to be understood literally. But he holds that,
since heaven is not any definite place, and since the meaning of the narrative
can only be that Jesus now left the spatial, sensible world, one must assume that
Jesus only raised himself a little way up into the air, which was to serve the
apostles as a sign that he was now assuming a wholly different kind of existence
from the sensible one! To this tasteless interpretation of the old account
Martensen is led by his claim that the Bible, by heaven, does not mean a place to
which one comes by long roads. He starts from the symbolic meaning of the
representation of heaven---and yet he will take the old narrative literally. What
was a possible standpoint for a medieval theologian was an impossible standpoint
in the nineteenth century, or ought to have been.

There is nothing else to do than to lay the whole weight on the symbolic meaning.
It is the only way in which the old accounts can come into their full right,
religiously, poetically, and historically. In every ``edifying'' application one
in reality takes this path. Only thus can one rightly understand what has stirred
in the human minds to whom these accounts are due. The ``natural places'' for the
religious events are to be found in the interior of souls, and only there. All
legends and myths are only testimonies to what has gripped minds in a certain age
and set the imagination in motion. And it is the task of historical research to
find out what outer events have been contributory to these spiritual occurrences,
and thereby to make possible a searching psychological understanding of them.
Historical research itself must again, as already remarked, start from the marks
of reality that we everywhere employ where we seek to orient ourselves in the
surrounding world. How far one can come along this path, and what religious
significance the results of the historical investigation of the origin and
development of religion may acquire, is a question by itself, into which I shall
not here enter. Personally I believe that the so-called liberal theology
cherishes too great expectations in this respect. --- I return to the question of
the influence of the thought of the immeasurability of existence upon the
religious consciousness.
It is by no means the case that the problem here arising should hold only for the
popular religions. Every life-view, whether it be religious in the narrower sense
or not, must be touched by it. Every human reflection upon the worth of life and
the conditions of life must meet the consideration that human life, with its
interests and ideals, its victories and its defeats, occupies only a vanishing
place in the world we know. We can no longer make the whole world the stage for
the drama of life. The stage upon which this drama unfolds is limited and
vanishing---and what significance and worth, then, has the whole of the rest of
the immeasurable universe? It is remarkable to see how little weight, as a rule,
has been laid upon this question, both in philosophy and in theology, especially
after the first impressions of the new conception had faded.

For Bruno it was a spiritual liberation when the narrow boundaries of the world
fell before his reflection. For his poetic-philosophical conception, existence
now stood as one great unity, everywhere ensouled by the same spirit and borne up
by the same forces. And in a similar way Jacob Böhme too conceived the matter,
once he had worked his way out of the grievous tribulation in which he was placed
when the infinity of nature had first dawned upon him in contrast to the
limitation of the human world. For Holberg it still stood thus, that astronomy
had now made possible a higher concept of God than the one held not only by the
common man but also by many theologians. But later on the question is seldom
discussed, and is preferably pushed aside. Positivists and Romantics here place
themselves at bottom in much the same way as the theologians. Comte's religion of
humanity introduced a symbolism in which the great universe stood as a serving
background for the globe upon which human life unfolded, and even purely
scientifically he declined all closer inquiry into the more distant parts of the
universe. Hegel (and among us Heiberg) mocked at the ``bad'' infinity that
extension in space presented. In Hegel's opinion men have overrated the
importance of the Copernican system. Although his speculation was to give us the
fundamental forms of the development of the \emph{World}-spirit, he nevertheless
maintained that the earth was the most perfect globe, because it was there that a
spiritual life was found. And Grundtvig found that there was no interest in
infinite space and the Copernican system, unless one started from the assumption
that sun, moon, and all the stars exist for the sake of men.

There is here, on all standpoints, a contrast between world-view and life-view
that may well invite reflection. Most inwardly, every serious life-view attaches
itself to the narrow circle within which human tasks and ideals hold; but it is
impossible that the gaze should not again and again slide out to the infinite
background. We may hold fast that in the vanishing part of existence we inhabit
there have developed forces and forms of life of a nobler kind than can be traced
in the rest of the universe; but we do not have at our disposal forms of thought
that could let these forces appear as those which determine the whole infinite
course of the world. Out of that which for us has the highest value we cannot
derive or understand the orbits of the globes and the surging of the world-
processes. There appears to us a gulf between the realm of our ideals and the
world of matter, a gulf across which no bridge can be built. One may close one's
eyes to this fact, but it is more manly to look it in the face and to shape one's
faith and one's hope in such a way that they can stand their test on the
battlefield of reality. We must have the courage to hold fast that the stream of
force and life which in its course has brought us all that we know of the great
and the beautiful will also henceforth be able to make its way through the hard
masses of reality, even though we can connect no meaning with the notion that
these masses should exist only for its sake. ---

The intuitability which is the strength of all popular religions holds not only
with respect to space, but also with respect to time, and in a quite
corresponding way. Yet it is not in all popular religions that the representation
of time and the representation of a historical development play an essential
role. Where these representations come forward, they point to a higher ethical
stage of development. It is now no longer merely the momentary and sensuous
welfare of individuals or of small groups of men that is the guiding interest in
the shaping of religious representations; it is the salvation of a whole people,
and finally of the whole human race, that is at stake. The religious problems
then concern the meaning of the whole of human history here on earth. The history
of religion shows us here a peculiar contrast. For at the point in development
where the urge arose to find meaning in the course of time and the march of
history, there was a significant intellectual tendency that could satisfy this
urge only by deadening it: the course of time had no worth at all, and it was
only a matter of stopping it, of saving one's soul out of it, of passing over
into a timeless being. This path development took in India, especially in
Buddhism, and Greek thought, for which the unchangeable had a greater worth than
the changeable, went in a kindred direction. That train of thought, on the other
hand, which came to exert the greatest influence upon European humanity and upon
the whole Muhammadan world, acknowledged the reality of the course of time, but
found the possibility of salvation in a great course of development that leads
peoples and generations to a final decision of the mighty world-struggle between
good and evil powers within and outside the human world. It seems to be from the
Old Persian religion that this thought especially went out into the world,
whether it arose within this religion or goes further back. All eschatology, all
doctrine of a development toward a last judgment and a subsequent world of glory
or of perdition, seems for a very essential part to be of Persian origin. It has
influenced later Judaism and through this Christianity, and it is said (according
to Goldzieher) to have directly influenced Muhammadanism from its very beginning.

Just as the world-drama, as regards space, got its limited stage in the ancient
world-picture, so, as regards time, it got its clear and definite place between
the beginning of the world and the conclusion of the world, and the temporal
distance between these two points, the beginning of the first act and the close
of the last act, was just as little infinite as the spatial distance between
heaven and earth. Here too popular religion has had its strength in intuitability
and in the limitation that accompanies intuitability. Here man saw himself, with
the whole spiritual and material weal and woe of himself and his fellows, as a
link in a great course of development, determined by the common past, perhaps
sighing under the torment of the present, but looking with confidence toward the
redemption that the future was to bring. Popular religion has here given life a
great horizon, made possible an understanding of the individual's fate and
striving from a higher point of view. And a significant advance in the historical
sense was thereby made. All this, without the frame of intuitability being burst.
It became possible to form clear representations of past and future, images in
which all the tension and longing, as well as all the jubilation and sense of
victory, that life in time can release, could find their strong, often exalted
and imperishable expression.

When reflection awakens and is carried through consistently, it fares, however,
with this intuitability as with the one that prevailed as regards space. The
fixed boundaries, the absolute beginning and the absolute end, fall as soon as
reflection applies the same method to the moments of beginning and conclusion as
to the intervening moments. In place of the short time (according to Parsism
12,000 years) that the world-drama was to last, there now enters the immeasurable
time, both as regards past and future. And the question then becomes how far
there can be given to that span of time which man's historical consciousness can
encompass such an absolute significance, such a rounded, measure-determined
character, as the world-conception of popular religion so naturally presented.
Can there now be any talk at all of a world-drama? The limitation in which great
art shows itself becomes impossible as regards the course of the world, and
therewith its dramatic character falls away.

As with the limitation of space, intuitability also presents, with the limitation
of time, the possibility of a feeling of at-homeness. It goes with the world as
with every single being: it has its origin, and it has its death. That the world
had grown old was a common belief in the first Christian centuries; among the
Jews too it was found, as is seen from the remarkable apocalyptic writing ``the
Fourth Book of Ezra.'' And one could then, amid the torments and disappointments
of the time, look forward with gladness to the approaching end. The near
conclusion was then no fantasy in which one now and again indulged; it stamped
the whole life-view, the whole ``Christianity of the New Testament.''

As with space, there also takes place, as regards time, a peculiar process of
inwardization in the religious consciousness itself. It seeks out beyond the
distances in time, as it seeks out beyond the distances in space. Not in a
distant future or past, but in present experience, will it have its object. The
ideal is not to have vanished or to be merely future, but is to make itself felt
in the very life that is lived in the now. Eternity is not to lie before or after
time, but is to be experienced in the midst of time---as psychological experience
too shows, that time disappears for us in those states that are filled with great
interests and ideas. In religious mysticism, whose germs are found in the
original religious documents themselves, but which was developed further in
various periods of the history of religion, the thought comes forward again and
again that the true eternity can be experienced in the midst of time, when the
will turns away from all that can scatter it and gathers itself about its highest
object. The religious consciousness, in such forms, approaches the conception
that alone becomes possible when the concept of time is thought through.

When reflection is carried through, we stand before the immeasurable time as
before the immeasurable space. We must in both respects find our place, the place
for our work and for the experiences we can make concerning the conditions of
life. And out from this place we then pursue the course of time forward and
backward, as far as we are able. But we must make ourselves familiar with the
fact that there lies time and events before the point with which we can begin, and
after the point with which we must end. If our interest is not limited to this
span of time, which is vanishing in relation to what the strict carrying-through
of the concept of time would require, then it must be a matter of holding fast the
conviction that what has real worth for a surveyable time need not lose it under
the continually continued course of time, even though this worth, through
manifold displacements, undergoes decisive changes. Within the periods of time we
can survey, we constantly see the phenomenon repeat itself, that what at first
stood only as means later acquired independent worth, and that what at first stood
as a final end later came to stand as means or preparation for more comprehensive
values. We must give up constructing the sequence of scenes in the world-drama; we
cannot decide in which act and in which scene our life falls; the very idea of a
drama can be applied only to a limited time. None of the dramas that art has
produced gives, moreover, by its conclusion any certainty as to how it goes with
the survivors later on. And thus are we precisely placed over against the course
of the world, when we let reflection prevail. Science is nothing other than
steady and patient reflection.

\medskip

2. With the expansion of the world-picture, as regards space and time, there
followed, in that remarkable age when modern science came into being, also the
demand for a stricter carrying-through of the concept of cause, and therewith a
new ideal for knowledge in comparison with the one that had prevailed in
antiquity and the Middle Ages. The principle was set up that, just as the events
or changes whose cause one sought to find lay before one as facts in experience,
so too the causes in which one found the explanation of them should be such as
could be pointed out in experience. Kepler, by whom this demand was especially
expressed with emphasis, had himself in his earlier writings spoken of causes
that did not satisfy this demand, and which he therefore later did not acknowledge
as ``\emph{valid causes}'' (\textit{veræ causæ}). He had thus taught that the
planets were held and carried in their orbits by star-spirits. But he became clear
that this explanation was scientifically invalid, since these spirits could not
themselves be pointed out in experience as acting in the way he had assumed. He
considered, on the other hand, magnetism (which he confused with gravity) a valid
cause of the planets' being held in their orbits, because it is in fact an
attracting force. However erroneous this conception was, it was nevertheless in
the spirit of the new research. Later Newton showed that the force which holds the
planet in its orbit follows the same law as the force which lets the flung stone
fall down to the earth. It was a great and exemplary instance of a ``valid cause.''
Just as every place in space must be determined in relation to other places, and
every moment in time in relation to other moments, so, by virtue of the principle
of valid causes, every event must be determined and explained by its relation to
other events that are just as demonstrable as itself. Hereby the world of
experience comes to stand as a coherent whole, and nature is interpreted by nature
itself, just as in a book one interprets a passage by means of other passages in
the same book. The weight lies here again on the method,---on the way in which one
asks, and the control one exercises upon the answer. It is a new organ that the
human spirit has here developed, and like every organ it can function only in its
own definite way. Thereby the tasks are set, and thereby the solutions are
conditioned. And it is in connection with this a new intellectual need that has
developed. Hitherto it was really only the formal sciences, logic and mathematics,
that had shaped themselves into independence; now the science of experience---in
its three principal forms: as natural science, psychology, and history---was in
principle liberated from theological and speculative points of view. And the
prospect over the new path that had been opened often awakened such enthusiasm
that men proclaimed in dogmatic fashion that everything could be explained by the
principle of valid causes, and that the causal series was absolutely coherent. Men
took the path of construction instead of waiting until observation and experiment
could show where the causes in the particular cases lay. It was a confusion of
result and task, of facts and principles, of system and method. But even when men
entered upon such transgressions, they were not mistaken in this, that they were
on the right path. It was the spirit of research that had become conscious of its
means and its powers, and therewith had also become conscious of its ideal. For it
was a great ideal that now presented itself to knowledge: to shape an ever-growing
number of experiences into an ever closer and more intimate unity-connection. It
is an ideal that at any given time can be satisfied only approximately; it
proclaims at every step an \emph{Excelsior!}---leads ever out to greater tasks.
Dogmatic, therefore, science cannot become, when it understands itself. Already
the immeasurability of space and time excludes a conclusion of our inquiry.
Something new can always come forward. And that not only beyond the boundaries of
the space and the time that at a given time one has been able to survey, but also
within the boundaries. New differences can turn up where hitherto there seemed to
be likeness or unbroken connection; an ever richer fullness of nuances and
transitions can be discovered within that which seemed to be so well known. The
world can constantly become greater for us, both in the great and in the small.
The very concept of world comes to stand as an ideal concept that cannot be
carried through, but constantly opens new horizons. And we may not even assert
that existence is comprehensible or understandable, measurable for our thought. It
might be compared with $\sqrt{2}$, the irrational number, which perhaps can be
determined progressively through the finding of ever more decimals, but cannot
receive any conclusive determination. The concept of world as a completed whole is
really not at all a scientific concept. It is a concept that the religious
consciousness has formed in its impatience, in order to get a sharp contrast to
its representation of God, but which for science denotes nothing given, only a
constant task. For science there is just as great difficulty with the concept of
world as with the concept of God.

When religion in its classical periods not only gave help and consolation in the
course of life, but also satisfied the urge to knowledge, it was not the principle
of the ``valid causes'' that it called upon one to apply. Its explanation of
nature was \emph{animistic}, that is, it maintained the intervention of personal
beings as ground of explanation. Just as the world as a whole was assumed to be
created by an act of will, so also the several events in nature, in the life of
the soul, and in history were assumed to be effected by the acting of personal
beings. Especially the natural events that intervene strongly in human destinies
were explained in this way. Storm and earthquake, as well as rain and fruitful
seasons, found their sufficient ground in the will of Providence, which now
punished or tried, now rewarded and blessed. Decisive changes in the domain of the
life of the soul were explained not by psychological causes, but by divine
intervention---as when the cause of Achilles' subduing his wrath toward Agamemnon
is explained by Athena's being sent down from Olympus to talk him into reason. The
march of history is interpreted by the prophets as the execution of a divine plan,
which bends the hearts of princes and turns the fortunes of peoples according to
its purposes. In a faded form the animistic explanation appears in the so-called
teleological mode of explanation, in which one speaks in a more indefinite way of
final causes, without being so consistent as to regard these ends as expressions
of personal wills.

Hardly ever, however, has the animistic explanation of nature been consistently
carried through. Already daily experience led to the assumption of a fixed
connection between events that both lay before one as facts, so that the one came
to stand as cause of the other. One let the incipient inquiry go its own ways a
longer or shorter stretch, but regarded the causes it found as derived and
subordinate; the real causes, the ``first causes,'' were personal wills. And one
assumed, moreover, that even in domains where as a rule the derived causes
operated, a higher cause could break through and---for the benefit of the divine
purposes---bring about a change that could not be explained by the ``natural''
laws.

There stand here, then, two systems over against each other, two explanations of
the world, the animistic and the naturalistic. There can really be no negotiation
between them, for they do not acknowledge each other's first principles. They ask
in different ways and therefore also get different answers. Over against the
miracle (which arises both where a ``first'' cause concludes the series of
``derived'' causes, and where it directly intervenes in the course of events)
science must declare that it cannot disprove it; but it adds that, according to its
principle and its method, it cannot catch sight of it at all. It can never come
face to face with a miracle; for where it cannot link one event to another
according to a definite law, it stands at a boundary, in any case a provisional
boundary. It then does not understand the event in question; it ends in a problem.
The two systems clash more and more in the human consciousness, and the victory
will belong to that one of them which can unfold itself most consistently and
serve as a binding tie for the greatest number of exactly determined experiences.
Could existence, as it actually lies before us, be better explained in an
animistic than in a naturalistic way, then the victory would be on the side of the
religious world-view. But what has been attained of real and abiding
understanding has not been won along that path.

It might seem as though an explanation of nature that sees in personal forces the
true causes must have its great justification and necessity already in this, that
man, the knowing being, is himself a personal being. Since all explanation
consists in leading the hitherto unknown back to the known, one might hold that it
must be the personal forces, which we immediately feel and know, that must furnish
the ground for all understanding of the rest of existence. This is now, in a
certain sense, also correct. We cannot leap over our own shadow, and the whole of
our knowledge of the world is conditioned by the human standpoint and human
nature. It is out from our own position that we orient ourselves in existence as
best we can. But the decisive thing here must be which sides and properties of our
nature we can make use of, when a valid and abiding understanding is to be won.
Experience teaches us, often through disappointment and pain, that it is not all
our purposes and expectations that we can get satisfied. It is originally a
practical interest that leads to putting definite questions to existence and
demanding definite answers of it. We seek means to attain our ends, and as soon as
we have discovered that we can attain our ends (those we can attain at all) only
along definite paths, by the application of quite definite means, then we already
stand at the causal relation that all science seeks after, the law that links one
event to another. Before such a law our wishes and hopes must bow; if they do not
fall away, they must often seek other outlets or more distant prospects. But
through the discovery of such a law a little insight has been won into the
constitution of existence, and the more such laws we can find, the better
understanding we shall have. It may then gradually be a motive other than the
original practical need that leads to the seeking of such laws. There can be a
quite immediate satisfaction in finding connection in the hitherto scattered,
whether such connection can point the way to the attainment of our ends or not.
Or rather, it can itself become one of our essential ends to win understanding in
the way it can actually be attained. It is precisely at this point that the
significant division of labour enters, which I mentioned at the beginning of this
exposition. It is, to be sure, a single side of the human personality's nature
that expresses itself in the intellectual need and capacity, as it appears in its
specially developed form. But what one must not overlook is that in this need and
capacity the very nature of personality lays itself open to the day. A personality
always presupposes an inner unity and connection between all the elements of the
life of the soul, all its forces and drives. And this holds not only for the
personal life considered as a whole, but also of every single endeavour to which
it devotes itself with full energy. Personality cannot in any single domain deny
its own peculiar essence. Within its knowledge it will therefore strive after as
great and unbroken a connection as possible. It will not cease to lay stone upon
stone in order to raise a building of knowledge that can stand as a counter-image
of the inner connection that makes itself known---in any case as an urge---in its
own essence. On this rests the deep personal satisfaction that scientific work can
yield, the passion it can awaken, and the power it can have to fill out a human
life. Personality and science are thus not foreign to each other, even though it
is not all elements of the soul that can in equal degree co-operate in scientific
work. Personality always spans more than its own intellectual work can encompass
and penetrate. Besides the task which both honesty and intellectual integrity set,
of procuring the most thorough orientation in the world that is possible from the
position on which the individual personality stands, there is set for it moreover
the task of working together its various interests and endeavours, especially of
letting the results of its intellectual work enter into as close a connection as
possible with the whole life of feeling and will, with all that is object of hope
and faith. That this great task is more difficult to solve than it was for other
cultural periods gives no right to knock something off it from the one side or the
other. We must run the intellectual line out according to its own laws, before it
can in harmonious fashion be woven into the whole life of personality. It is the
great task that is ever more definitely set for the people of our time---for those
of them who do not dream with open eyes.
The path upon which the science of experience has thus entered is therefore no
accident. The empirical-scientific method has been wrung from human thought by
existence itself, which will let itself be questioned, and itself answer
questions, only in this way. More and more, to be sure, it becomes customary to
speak of working methods and working hypotheses, and, as we shall see later, it is
of great importance not to confuse the paths along which we must walk in order to
catch sight of the truth with the inmost essence of existence itself. But a wholly
external and indifferent relation between the path and the goal is nevertheless
unthinkable. It serves to characterize existence that its laws can be found only
along these definite paths. Indirectly one comes to know the nature of existence
through the working hypotheses themselves.

But seen from another side too, the procedure of science has its good ground,
namely if it is correct, what I sought above to show, that there is a deep analogy
between this procedure and the very essence of personality. It is naturally not
only this that has occasioned the method of more recent science; but it can serve
to explain the deep conviction with which, again and again, despite many
disappointments and mistakes, men return to this method and seek to develop it
further and to apply it to new domains.

In the last generation there has happened in science various things which, to an
external view, might seem to shake the reliability of the points of view that in
the foregoing age were so securely relied upon. Men have accordingly been busy
proclaiming a bankruptcy. They overlooked what the whole history of science
shows, that in a certain sense the problems become more numerous the further our
knowledge advances. Sharp problems exist at all only for those who think and
inquire. And the more various experiences we make, the more tasks are set; for it
is then a matter of bringing the new experiences into connection with one another
and with the earlier experiences, which cannot always happen without alteration of
earlier assumptions.

With the discovery of radium the principle of the conservation of energy seemed
thus to be shaken. Here there showed itself a constant radiation of energy,
without its being possible to point out how this energy had been deposited in this
substance. It then set new questions and gave occasion for new reflection. But it
changes nothing in the old method, as regards the matter of principle. --- At
about the same time it was shown in biology that organisms can undergo sudden
changes, whereby new types are formed. Thereby it became necessary to give up the
assumption put forward by Darwin, that new types arise by small variations being
favoured and developed under the influence of the struggle for existence. The
mutations, the sudden changes of type, occur namely independently of the struggle
for existence; only after the mutation has taken place does the question arise
whether the new type can survive. But here too the presupposition is not given up,
that every event must hang together with another event. The task becomes to find
the inner and outer conditions under which the mutation takes place. --- A third
example we may take from psychology, where in the last decades attention has from
various sides been drawn to sudden changes of character and personality, new
formations in the psychic domain, especially within the life of feeling and will.
Here too the conditions show themselves more complicated than one had perhaps for
a time believed. But the points of view in principle continue to be the same;
only there is greater difficulty in carrying them through.

The conditions indicated here are a wholesome warning against dogmatism in
science. We have again been reminded that there is more in heaven and earth
than we thought of in our philosophy. But truly scientific thinking does not need
such a warning; it obeys Hamlet's counsel: to bid the new and wondrous welcome,
even though our philosophy had not previously thought of it\footnote{\emph{Hamlet}
I,~5. (Horatio.) O day and night, but this is wondrous strange!\\ (Hamlet.)
\emph{And therefore as a stranger give it welcome.}\\ There are more things in
heaven and earth, Horatio,\\ than are dreamt of in your philosophy.}. --- It is
precisely thought and science themselves that have opened the great horizons which
at the same time set us the great problems. It does not fall to those whose faith
rests essentially on old traditions, which they do not always confirm by their own
life-experience, to come forward with admonitions to keep the gaze open. ---

A special ground for holding fast the principle of the ``valid causes'' lies in
this, that only then does it become possible to make trial of, and get
confirmation of, the cause one believes oneself to have found. When the cause of
an event or a change we have experienced is sought in a preceding event that has
itself appeared in experience, then there is the possibility of calling forth this
preceding event again and seeing whether the presumed effect then also enters
anew. If, on the other hand, one appeals to causes that cannot themselves be
experienced, directly or indirectly, it is not possible to get such a
confirmation. A miracle does not let itself be repeated at will. It cannot be
shown whether the personal forces or powers assumed as its cause are present anew,
and present in quite the same way as the first time. Here lies the real ground of
animism's intellectual barrenness. It can give expression to the wonder at the
striking, the unexpected, but it cannot in the right way bid ``the stranger''
welcome. For to that belongs precisely this, that the new becomes, through
repetition, familiar to us and is taken up into the circle of our thoughts, even
though this can perhaps be done only after a change in our earlier way of thinking.
Wonder is only the beginning of wisdom.

In the trial of causal explanation that is made by calling forth the presumed
cause anew, the presupposition is that the very event that is to be explained has
really taken place, is not a mere illusion. Wonder can often lead to a hasty
assumption of reality. But here too the importance of the principle of the
``valid causes'' shows itself. For when we ask what mark we have of an event's
reality, it appears that we are here always referred to testing whether it accords
with the experiences we have previously made, fits into the whole circle of events
and relations that already stand fast for us. It is the mark we quite
involuntarily employ when we are in doubt as to whether something is real or not.
We investigate it with our various senses, and we compare the present testimony of
the senses with their earlier statements. We use the experiences of others, which
we believe we can rely upon. And our doubt falls away only when this whole circle
of experiences and recollections forms a coherent whole, without
self-contradictions and without leaps. We apply, then, one and the same method
when we test the reality of the single event, and when we seek the cause of the
event thus tested. Historical criticism therefore also necessarily applies this
method, whatever accounts it may stand before. It investigates how far the several
testimonies agree,---whether the one is more than a mere repetition of the
other,---whether what is reported accords with what experience teaches us about
the course of events in the material and in the spiritual domain,---whether the
content of the account cannot find its explanation in the narrator's personality
and state of mind or the traditions in which he has grown up, and so on. The mark
of reality is never anything single. Reality we discover by a comprehending
activity of thought. A single beholding, a momentary vision, however much it may
grip, contains no certainty of truth. What thus suddenly grips us can become an
incitement to inquiry, to all-sided testing, and it may then show itself that it
was a discovery that was made. But it contains no probative force, neither as
regards the several objects each for itself, nor as regards the explanation of
them. When now a critic stands before an account that contains a miracle, the
question for him will be whether it is more probable that such an event has
happened, or that the narrator has been mistaken for one reason or another. Here
the question will be, among other things, whether the account and its content
cannot be explained out of the spiritual conditions of the time. An account is
itself an event, a fact that lies before one; it is not only what is reported, but
also the fact that it is reported, that can set investigation going. Historical
criticism therefore applies, in its way, quite the same method that science
applies in all domains.

Now there is, to be sure, a significant difference between the demonstration of
the validity of general laws and the adducing of proof for the reality of a single
definite historical event, namely the difference that a law must be able to
demonstrate itself ever anew, whereas a historical event enters only once. With
this it hangs together that it is on the whole more difficult to show the validity
of the causal principle, the more special and concrete an object is. There is,
with respect to the possibility of repetition, and therewith also the possibility
of scientific confirmation, a whole scale of degrees from the simplest to the most
complicated objects. In mechanics, physics, and chemistry it is comparatively easy
to bring about repetition, to come to stand before quite the same conditions as on
earlier occasions and then test whether what follows also shows itself to be the
same as then.
In the organic domain it is far more difficult; the state of an organism or an
organ varies far more, depends on far more various conditions, than a mechanical
motion or a chemical state. States of the soul vibrate and shift still more, and
historical events, in which forces of nature and manifold organic and psychic
conditions work together, show a minimum of repetition. Yet it is only differences
of degree that in this respect make themselves felt between historical events and
natural events. Not even in physical nature is there absolute repetition. Where we
assume such, it rests on inexact apprehension; we do not get all the nuances and
conditions with us. Therefore the method of investigating what really lies before
one is in principle the same everywhere. Historical criticism has only to do with a
larger and more complicated aggregate of relations and circumstances. The same
holds of historical characterization, the final determination of a historical
event or of a historical personality's peculiarity and significance. But this side
of the matter I shall come to touch upon later on. ---

We return to the miracle, and to animism as a whole, in order to call attention to
the fact that within the religious consciousness itself a criticism can make itself
felt of the externality and doubleness that animistic explanation, and especially
the assumption of miracles, presupposes or leads to. The personal forces here
appealed to are of a wholly different kind from that which they are to explain.
They come from without and set the thing in operation. There is then no inner
kinship between cause and effect. The divine beings come to stand as those who
operate in a world that is foreign to them. We find here the earlier-mentioned urge
in the religious consciousness to raise its object as high as possible above all
that is finite and earthly. But when this tendency has led to there arising an
external and distant relation between the deity and the world within which man
himself has his place, there is a reaction in the opposite direction. Inwardness
demands its right. The deity is to be near us, to work in us ``from within,'' not
``from without''---and then it cannot stand in a wholly external relation to ``the
world'' and to the world-forces either. From the religious side too one cites the
poet's words: ``What sort of a God would it be, who could only strike upon the
world from without?'' If the world were something that had independence in relation
to the deity, it would itself be a god. God would have to use some of his power to
work upon it. In particular, the miracle would really be a testimony to a struggle
between two gods and could in any case be thought only in a polytheistic religion.
For the most inward religiosity the external distinctions fall away, and all the
externality that is reported of the deity's intervention acquires in the end only
figurative significance. This movement is continued in mysticism, but there is a
mystical element in every inward religious consciousness. And this tendency seems,
so far as one can judge, gradually to become of such preponderant importance that
the miracles, which formerly were the most important means of proof for religious
faith, now stand rather as hindrances---hindrances one quietly wishes away, if not
on scientific, then precisely on religious grounds. In every edifying application
of miracle-narratives symbolism is accordingly clearly prominent. As with space
and with time, so here too, with causal explanation, there goes a tendency to
symbolizing through later stages of religious consciousness. What is lacking is
often precisely only this, that one becomes clearly conscious of the symbolic
character.

Descartes declared in his time that in his explanation of nature he would not
apply the so-called final causes, that is, explain the coming-to-be and
constitution of things from the ends they are supposed to serve. It was in his
eyes presumptuous, because we thereby ascribed to ourselves an acquaintance with
the intention God had had in bringing forth things. It is indeed clear that such a
knowledge of the ends of things is presupposed in teleological explanation. And it
is presupposed still more clearly in the assumption of a miracle. The presupposition
of such an assumption is, after all, that there has been sufficient occasion for
the higher powers to intervene in the customary course of things. And how can that
be known? Pascal held that such an occasion was at hand when a young girl was
healed in the church of Port-Royal; but the opponents of Port-Royal denied it.
Only a revelation could decide the matter, and it is therefore consistent that
already the medieval Church reserved to itself the right to decide whether in the
particular case a miracle was present or not. In the most recent time the Catholic
Church has inculcated this anew on the occasion of the many apparitions of the
Madonna in the last half-century.

But it is not only the divine ends that must be known in order that one may assume
a miracle. It must also be possible to decide definitely that a deviation from the
law-governed order of nature has really occurred, and this really presupposes
omniscience too, since it presupposes insight into the limits of nature's
conformity to law, limits which science has not been able to find---and never will
be able to find. ---

The animistic explanation of nature, and especially the belief in miracles, is due
in part to a natural urge to understand and explain, an urge that at earlier stages
of development could not be satisfied by other means of thought. But the main thing
was, and has more and more become, that such an explanation of nature made it
easier to conceive existence and the course of things in existence as determined by
that which for the human spirit stood as the highest good. Personal beings act
according to ends, strive after goods for themselves and others. When it is now
such beings to whose intervention the great changes in nature, in the minds of men,
and in history are due, then in the end a thought of purpose stands as the decisive
thing, and when man can believe that his own highest welfare is comprised in this
thought, then he can feel secure amid the vicissitudes of fate. What stands for him
as the highest of value is then also the last cause of what goes on in the world of
reality.

But at bottom the religious view has then, by the help of animism, reinterpreted
reality. The strict mark of reality, which, as shown above, consists always in the
possibility of linking event to event according to the principles of the ``valid
causes,'' is disregarded, and the great problem of the relation between the ideal
realm of values and real actuality is therefore not seriously posed. One does not
look the matter sharply in the face, hastens with too great haste to find
consolation, letting one's wishes determine what is to be truth and reality.
Whatever the religion of the future may become, and whether or not the word
religion will continue to fit the faith or the life-view that will develop, it will
certainly in any case become its distinctive mark that it poses the great questions
of life as sharply and definitely as can first become possible when the great
spiritual work that science performs becomes seriously co-determining both in the
questions and in the control over the answers.

In the God-concept of the higher popular religions two thoughts are fused into one:
the idea of the highest value and the idea of the absolute cause. Such a unity
scientific thinking cannot reach. It finds in experience a shifting content between
the valuable and the real. The Good, the True, and the Beautiful (to use this old
trinity) must constantly struggle to survive and must often wrest from reality
their right to be. Out from these ideas we cannot understand the whole of
existence. We can perhaps find and follow a stream of light and harmony and hope
that it will always be able to continue its course. But this continued course
itself we cannot construct. What at a given time stands as the highest of value,
the true Good, True, and Beautiful, is for a later time often only of passing
significance, is a means or at least a transitional form to other values. There
constantly takes place, in the great as in the small, displacements and
transitions that make it impossible to say which values will continue to assert
themselves. A highest value is therefore a concept that we cannot determine once
and for all. The series of values we cannot think as concluded. And as it goes with
this series, so it goes also with the series of causes. It constantly lengthens
itself for us, and the new, future links in the series we cannot in advance derive
from the preceding, just as little as we can lead the whole series back to an
absolute beginning-link. As far as we see forward or backward in both series, there
constantly present themselves new tasks and problems. To think both series brought
to a conclusion each for itself, and moreover the one led back to the other, is to
want to find the philosophers' stone. And if one begins with the belief that it has
been found, then naturally all riddles become easy---all too easy---to interpret.
One might ask whether scientific thinking does not then need a conclusion,---whether
it cannot form an idea that denotes the last, comprehensive point of view. To this
it must be answered that such a conclusion can in any case not be reached by
following the series of causes forward or backward; that leads only to ever new
tasks, ever new work. Not as a link in the series, whether as first or last link,
as beginning or as end, can thought find a concluding idea. But the very fact that
there is a causal series, a world-order that reveals itself in the conformity to
law, is for thought a characteristic mark of existence, such as we know it. Therein
makes itself known a great unity, a connection in the great as in the small, not
only in the outer world of matter, but in our own inner being, surely further
within than we ourselves can penetrate. This idea is the scientific limit-concept,
which cannot receive a positive form except by the help of more or less poetic
analogies. We are here cut off from employing the only procedure we have for
finding positive determinations of concepts, namely comparison with other objects.
That which is to be the bond between all our experiences cannot itself be bound to
anything else. ``That which holds the world together in its inmost core,''---to use
Faust's expression,---cannot find its expression in any of our concepts, because
all our concepts are drawn from more special and limited domains. And there is
added to this, as will later be more closely developed, that the significance of
this limit-concept lies in the understanding that it is possible by its help to
reach of the several individual existences that experience presents. The
``concluding'' idea contains a direction to constant work. It is set up not as an
object of contemplation, but as the guiding thought whose leadership we have become
conscious of. The ``concluding'' idea therefore really brings no conclusion, but
leads us back again to our workplace in the world of experience.

Every personal life-view will now be conditioned not only by the purely
intellectual motives that lead to the idea of a collection of causes, a great
world-order, but especially by the place which that which stands for the
personality as the highest value occupies within this world-order. It will then
depend on what experiences the individual has made with respect to the struggle of
values to survive. The individually different character that faith or the life-view
acquires in the various personalities will depend partly on the values whose fate
is object of hope and fear, partly on the experiences that have been made, or are
believed to have been made, with respect to this their fate. --- Upon this whole
question I shall not here enter more closely. It concerns comparative religious
psychology. A personal life-view must by the nature of the case have a far more
special and concrete character than a scientific theory has. It arises through the
combination of the scientific knowledge the individual has acquired with practical
experiences of life. --- Some remarks on how the relation between faith and
scientific thinking shapes itself, or must shape itself, in our time shall
nevertheless here be set forth.

\medskip

3. Just as in nature various forces are at work, between which the natural
scientist seeks to find a connection in order if possible to demonstrate their
unity, so too in the human soul various, often contending forces make themselves
felt, and it is a task, not only for the man of science, but for everyone, to find
a connection, and if possible a harmony, among them.

Thinking and faith are two such forces or tendencies, often influencing each other,
often in sharp strife, not only because some individuals will hold to thinking
alone, others to faith alone, but also because they can appear as opposites in the
individual's inner being and produce a cleavage there. And yet they both spring, in
their inmost ground, from the same source: the human personality.

Thought seeks order and law-governed connection among the observations, and where
it can find no order and law, it seeks to widen experience, to make new
observations, led by the hope that what, so long as it stood isolated, seemed
accidental and capricious, may show itself to stand as a link in a whole
connection, when it is held together with a greater quantity of experiences. Often
it must, however, halt, because no solution can be found; the observations are
either too few, or they are in their manifoldness and diversity so inexhaustible
that no common point of view can be set up. Therefore thought does not on that
account give up its demands and its ideals. An absolute conclusion cannot be
reached; but then we end with a question, in a problem, and to that extent wonder
can be not only the beginning of wisdom, but also the mood that corresponds to our
wisdom's last word. Besides, there is, as we have seen, an idea that accompanies
thought on its way from first to last, because it both denotes the constant task
that thought sets itself, and is at the same time an expression of what the work of
thought leads to when it succeeds: the idea of a great world-order that makes
itself known in the laws of events and in the connection among our experiences.
Thought works with a high horizon, even though this horizon at the same time always
denotes the boundary of what it can reach.

While thought thus strives outward, goes from problem to problem, constantly seeks
after new material and after an extension of its horizon, faith denotes a gathering,
a concentration of mind and will about a single representation, a single object,
which for the individual casts light over life, because he finds expressed in it
what gives life value, and at the same time what secures the survival of this value
through the surging of fate.

The relation between thought and faith is like the relation between a centrifugal
and a centripetal force. The highest thing would be to unite the mobility of
thought and its great horizon with the inwardness and concentration of faith. That
is also the task that is set for every serious striving. But both in the history of
the race and in the life of the individual it shows itself that this harmony is not
so easy to attain. There arises again and again a strife between the opposed forces.
Perhaps it is an eternal strife. As long as life lasts, there come new observations
that thought is to order and comprehend, and fate brings new experiences of life to
which faith is to take up a standpoint. The task of reaching a harmony between
thought and faith is therefore set again and again under new conditions. The strife
will then constantly be carried on under new forms.

In the nineteenth century it has thus run through various stages. At the beginning
of the century Romanticism announced that the spiritual division of labour was to
be over. Spiritual life was to constitute a unity, which under various forms
appeared in art, in science, and in religion. The unhappy strife was to be over.
What thought comprehended was to be only another expression of what art beheld and
what faith surmised. In Denmark it was Henrik Steffens who proclaimed a return to a
spiritual condition such as it was supposed to have been present in a bygone time
that did not know our modern oppositions and distinctions,---a time that is depicted
in the following way: ``In living action was impressed all that we now seek by
words, all the revolutions and developments of nature, all the mysteries of
religion. Life itself, fresh as it sprang from the deity, had immediately the
significance that we now only surmise. Therefore all representations were living. No
trace of the dividing death which in more recent times has made the sciences and
poetry, and through them religion, so monstrous.'' One sees that Steffens derives
the difficulties in the religious domain from the emancipation and specialization of
science and art. The train of thought he asserted played a great role in Danish
spiritual life, from Steffens' Lectures (1802) to Martensen's Dogmatics (1849). But
our two greatest religious personalities, Grundtvig and Kierkegaard, protested
energetically against it, and their appearance has led to a sharpening of the
opposition between the two spiritual forces. To this was added, in the last part of
the century, the mighty upsurge of natural science, the pressing-forward of the
historical method in all spiritual domains, and the reawakening critical
philosophy's testing of the validity of our representations---also of the validity
of those representations in which faith procures itself expression.
In the face of this state of things it might lie near to take refuge in an absolute
cleavage between thought and faith, so that the task became to hold them apart from
each other, since this was the only way in which both could be held fast. Under the
influence of Kierkegaard's writings and Rasmus Nielsen's lectures I tried this way
out in my youth, but both life-experience and reflection taught me that it could
not lead to the goal. The two spiritual powers cannot stand wholly indifferent
toward each other. Thought can make faith itself, its motives, its ideals, the
representations in which it finds expression, the object of its inquiry, and on the
other hand the work of thought itself springs from an urge that is akin to the one
that leads to faith. There is then reason to institute a closer investigation of
their mutual relation, in order thereby to recognize under what conditions the
collision between them, which cannot be avoided, must be supposed likely to take
place in the future. We look somewhat differently upon the whole question now than
Romanticism and Positivism, the two most important tendencies of thought in the
nineteenth century, did.

\medskip

On two different points, both of which have been briefly touched upon in the
foregoing, there shows itself an inner kinship between scientific thinking and
personal life.

a. All science is a striving to find connection among scattered and heterogeneous
observations. At all times men have felt an urge to this. Involuntarily, in the
human spirit, the many observations are transformed into coherent images and
thoughts, just as involuntarily as the nutritive substances in the organism are
transformed into flesh and blood. The difference between the various stages of
development of thinking rests on how many and how exact observations one has at
one's disposal, and how carefully one compares and connects them. That there must
be a connection is a presupposition that involuntarily makes itself felt,---in what
for us stands as childish superstition, as well as in what we call science, and
which perhaps in a thousand years one will not call so. In Egyptian mythology the
sun that rose in the morning was not the same as the one that had set in the
evening; but they were nevertheless put in connection, the morning sun being
assumed to be the son of the evening sun. The falling of the stone to the earth and
the movements of the moon about the earth were found by Newton to be examples of the
same law, and Ørsted found an inner connection between electricity and magnetism,
which had hitherto stood as different forces. The passage-graves of the North and
the domed tombs of Greek antiquity have been shown by the archaeologists in their
historical connection. Psychology shows that the coming and going of our
representations are subject to definite conditions, which can be formulated in
general outline.

The urge to, and the joy in, finding connection is rooted, as we have already seen,
deep in the very nature of personality. We call a man a personality because there
makes itself known a connection among his endeavours, his moods, and his
representations, and the more energetic, peculiar, and self-acquired this connection
is, in the higher degree do we ascribe personality to him. Never is a personality
exhausted in a single state or a single utterance, however peculiar this may be. The
personal life has at once a fullness and a unity that points out beyond each single
state, and the measure of personality consists partly in the compass that is
spanned, partly in the inwardness with which it is comprehended. Science is now
nothing other than one of the forms under which the urge to fullness and unity
expresses itself, a form that in intellectual natures and scientific geniuses
appears in a marked degree. It is always itself a form of personal life and is no
accident in the life of the spirit, and can therefore never come to stand in a
purely external relation to it. For some natures there can be bound up a great
spiritual pain with having to halt at incoherent or perhaps even contradictory
observations and representations, just as for every developed personality there is
pain in experiences, promptings, or representations that cannot be brought into
harmonious relation with the character and with the life-view as a whole. Intellectual
integrity, which with all its ability and power gathers the observations that can be
made and connects them in as clear a connection of thought as possible, is then also
a virtue that acquires a greater significance and necessity the more the kinship
between science and personality becomes clear. Not everyone can be a man of science,
but along ever more paths it becomes possible to enter into relation with the
treasure of knowledge that the race has won on its way, and the individual stands
answerable for the way in which such a possibility is used by him. It is a form of
love of truth on which in earlier times, for good reasons, so much weight could not
be laid. The love of truth could then appear only as honesty, care for agreement
between confession and faith, or as personal truth, care for agreement between faith
and the personality's real inner life and striving. Already the circumstance that
there comes to be room for a new kind of love of truth is a testimony that the
opposition between thinking and faith cannot remain standing as absolute.
\medskip

b. One will perhaps say that science nevertheless leads away from personal life, in
that it seeks general laws, over against which the several cases, events, and
individuals have only vanishing significance. Far out beyond the individual
destinies and the narrow circles within which they unfold, out into the
immeasurable worlds of space and time, within which the single human life is not to
be descried, thought seeks in its zeal to trace out the order and conformity to law
of existence. Every personality is precisely single and individual, bound to a
definite time and place, but science seeks the universal, the general. It is perhaps
at this point that one finds the sharpest opposition between science and personality.

It is true that science has celebrated its highest triumphs by demonstrating general
laws, and this will provisionally also continue to be its most important task. But it
would not be justified to regard this as science's highest task and final ideal.
There will always remain the task of showing how the several existences find their
place within the great world-order, which they at once are borne by and themselves
help to sustain. Insight into general laws is in the end a means to understanding the
several, individual existences. In order to find the general laws, thought must
resolve these existences into their elements; for the most comprehensive laws hold
for the elements of things in their mutual relations. Things must be made simpler
than they are in reality, in order that the most general points of view may be
applied. But when this work is done, it is again a matter of connecting what one has
separated and of letting the general laws and points of view co-operate toward an
understanding of the several, definite wholes that experience presents. The role
that developmental hypotheses have played in the last hundred and fifty years points
in this direction. The general laws that have been found are used to form a
conception of the solar system, of the development of the globes and of living
beings. Thereby natural science itself acquires a historical character. In psychology
it is a matter not merely of finding general laws for the interplay of
representations and drives, but of using them toward an understanding of definite
individual personalities. And social statistics is to help us toward an
understanding of the development of society. It is the interplay of laws in the real
existences that in the end matters. Every individual is, as H. C. Ørsted said, a
compound unity; its subsistence is due to a co-operation of laws. And as Leibniz
still earlier expressed it: an individuality we come to know through the law of
development that holds for its various states. Science's task is in the end
biographical. The complete understanding of the wholes that experience shows us in
their constant collision and struggle to subsist,---star-systems, globes, organisms,
personalities, societies,---would be science's real conclusion, but it will
constantly stand as an ideal. There presents itself, notably in all developmental
doctrine, the already-mentioned circumstance that it is difficult, if not impossible,
to make trial of the correctness of the hypotheses, since we here stand at last
before \textit{unica}, before existences that appear only once in our experience.
There can then here always be talk only of approximation. There will always be an
irrational relation between the general and the individual.

In connection with this stands the fact that one more and more regards scientific
principles and laws as working hypotheses, as presuppositions whose significance
consists solely in their setting thought to work. When it is taught, for example,
that in nature no matter and no force arises or perishes, then the significance of
this proposition rests on its calling upon us, every time matter or force seems to
arise or perish, to seek what other matter or forces they arise from or are
converted into. Could other hypotheses set research in motion in a more fruitful
way, they would be preferable. The so-called pragmatic tendency in the theory of
knowledge has in our days put this view on its point, so that it can often look as
though the thoughts we form about existence really had nothing to do with existence
itself, and at bottom nothing to do with our own nature either, but were accidental
means of orientation that, whenever it should be, could be exchanged for others more
suitable. But, as already remarked earlier, it serves to characterize existence that
it can be understood only along certain definite paths, even though these paths are
not the same at all times. And we must add that it also serves to characterize the
human spirit that it can reach knowledge only by forming its thoughts in a certain
definite way. Long before Pragmatism the critical philosophy had inculcated that we
must not confuse the forms in which our thoughts move with the properties of
existence itself. Both tendencies, Criticism as well as Pragmatism, are symptoms
that science is becoming less and less dogmatic. They both warn, each from its point
of view, against regarding the views we employ to orient ourselves in existence as
``eternal truths.'' The truth of scientific propositions consists not in their
depicting existence for us without further ado, but in their containing directions
for fruitful spiritual work. Perhaps at some time wholly other principles and laws
than the present ones will be necessary in order to attain the fullness of
observations and the connection among them which science constantly seeks. We have
not exhausted the possibilities of existence and do not know all its demands, and we
have surely not exhausted the possibilities of forms of thought that lie in the human
spirit either.

This last consideration is not only of importance for the relation between thinking
and faith, but is of special importance in a time when so much is being done for the
popularization of science. The art of popularizing is so difficult because it is a
matter not merely of communicating finished results, dogmatically formed
propositions, but of giving a conception of the mode of work, of foundation and
limitation. The task must be to lead into the researcher's inner workshop, and that
can be done only indirectly, through intimations and by awakening self-thinking. It
is the old Socratic method that must be applied, if popularization is to mean
anything other than filling with finished representations whose origin and range are
not understood. ---
What holds for science, that a sharp distinction must be made between the objects
that lie before one and the interpretation we give them by means of our forms of
thought, our principles and hypotheses, holds also for all faith, all personal
life-view. If science has been too dogmatic, has been inclined to regard its
thoughts and hypotheses as eternal truths, then dogmatism has been able to spread
itself in a still higher degree in the world of faith. This is due to the
circumstance that tradition and habit play a greater role in the domain of faith
than in that of science. In science too, accustomed views and traditional
conceptions can long exercise their power over thought; consciously or unconsciously
the established exercises its influence. The long dominion of schools, tendencies,
and methods shows this. But tradition has a still greater power in the domain of
religion. Here it is perhaps even regarded as a necessity. It is not enough that a
solution of life's riddles is given, or that supports for hope and courage are
given; it must also stand fast where they are to be found, and the religious
consciousness feels really secure only when they are guaranteed through venerable
tradition. Authority has played and constantly plays a decisive role in the
religious domain, and the very concept of faith is as a rule understood not as an
expression of what the individual's life-experience and reflection have fastened in
his conviction, but signifies trust in, and adherence to, a handed-down
interpretation of life and its destinies. The gathering and concentration which,
psychologically seen, is the peculiar thing about faith, in contrast to the often
errant wandering of thought, appears most often as an adherence to a proclamation, a
historical model,---an adherence that has now the character of childlike trust, now
of obedience, now of passionate clinging. And even where the religious consciousness
gives itself over to its visions or promptings, immerses itself in its experiences of
life's worth and destiny, even then the handed-down forms and images will
involuntarily and unconsciously push themselves forward and stand as the only natural
expression of what has been seen and lived through. When a religion is in its first
period of development, there is a certain freedom for personal experience and
interpretation. But gradually, as the organization in dogma and cult advances, ever
narrower limits are set for the free spiritual life of individuals. Just as the
Church reserves to itself the decision of what is a miracle, so it also arrogates to
itself the decision of what is religious experience, and how it is to be interpreted.
The mystics of the Middle Ages therefore revised, in moments of composure, what had
been written down after their visions in the moments of enraptured ecstasy, so that
nothing should occur in it that ``ran counter to the utterances of the holy
Fathers.'' But involuntarily the handed-down and sanctioned representations were
woven already into the visions themselves. Where authority and tradition thus,
consciously and unconsciously, hold sway in the religious domain, it stands fast once
and for all what spiritual nourishment a man needs, what he needs in order to believe
and hope, in order to lead a noble life and hold himself up under life's struggle. No
heed is paid to the great differences among men themselves, which bring it about that
they cannot have quite the same models or quite the same support. And no heed is paid
to the fact that a personality is not finished at any stage or in any form; there
constantly come new experiences in which the won conviction is to stand its test, new
situations in which the life-yield hitherto attained is to prove its worth.
Personality can grow only with its greater aims, when it does not, without testing,
close off in definite views. Here is a point where the history of science contains an
instruction for personal life as a whole, in that it shows the necessity of raising
oneself above the accustomed and the near-at-hand and of adopting new points of view
in order to understand the new. Science can work with its highest thoughts without
letting them crystallize into fixed forms for all eternity. And it holds of all our
working thoughts as a whole that they exist not merely for their own sake, but for
the sake of life. This must then hold quite especially of the religious
representations. Just as the causal principle itself is an anticipation, an utterance
of an intellectual hope for which the work of thought seeks confirmation, so too
religious propositions will be able to preserve their significance, and perhaps first
rightly acquire their significance, when they are essentially understood as
expressions of the great hope for the preservation of the values of the spiritual
life, even though these values will assume new shapes under new conditions,---a hope
that is in reality the conscious or unconscious presupposition of all serious
striving.

There is nothing restless or Tantalean about this way of viewing things. It is once
and for all so, that neither in the inner nor in the outer world is there absolute
rest. At the summits of our striving, admiration for the life within us and around
us will naturally break forth and procure itself expression in the greatest images
that our thought and imagination command. But if it is not to become fruitless
raving, it must show itself active in continued striving, even though in a more
hidden way. There is a natural rhythm between the great enthusiasm and the quiet
work, a rhythm that holds in equal degree for the life of faith and for the life of
thought.
\medskip

4. As we have seen earlier, ever new tasks present themselves for science, and
existence stands as inexhaustible for it, despite its constant advance in relation
to its beginning stages. There is an irrational relation between thought and
existence. But a corresponding relation obtains between the personal human life and
the forms in which it expresses its faith and its hope. Since personality comprises
other elements than the intellectual, it will hardly ever be able to find in science
the full expression of what stirs within it. There will therefore---it is a
psychological assertion which closer investigation will constantly find confirmed in
the particular cases---appear at every standpoint a difference between knowledge and
faith. Every single personality is a centre for experiences that could be made
nowhere else in the world. The deeper the experience of life goes, the more
individual it is, and the more difficult becomes its conversion into general laws and
points of view; indeed, the mere communication in words already presents
difficulties. Only recently have the psychologists become seriously attentive to the
significance of the individual differences. Human life is like a meadow. Most people
see only the uniform green; one must bend the blades aside in order to become
attentive to the manifold life that stirs beneath them. Every man who is more than an
echo of his surroundings builds involuntarily his conscience and his faith upon the
experiences he has himself made. What gives life value, and what therefore determines
our greatest hope, everyone must learn on his own account, and all instruction can
only aim at helping the individual to make his experiences of life in his own way, as
the peculiar life-centre that he is.

This point has long since been emphasized by Danish thinkers. A. S. Ørsted declared,
in agreement with Treschow, that the peculiar and the individual cannot be exhausted
in concepts, and he draws from this, as regards ethics, the consequence that no
general moral commandments can exhaust what is required in the particular case, and
especially that such general rules cannot let the individual's peculiarity come into
its own. Everyone must therefore in a certain sense have his own ethics. Sibbern laid
great weight on the fact that development in all domains begins from scattered points,
so that the connection and the community forms the conclusion, not the beginning.
Thus spiritual development too begins in the several personalities, in which life
stirs in different ways, and only through the interaction among them is it then to
show itself whether community can come about. All that is later understood as
objective doctrine has had its first origin in the inmost of single souls. Especially
in the religious domain, unity has value only when it comes as a natural consequence
of the same knowledge being won by all. When one dogmatizes, one tears the truth
loose from its living root; one makes a mummy of it and lets it be handed down in
this shape from generation to generation. Poul Møller inculcated the concept of
personal truth in contrast to the affectation that lets a man believe and behave as
though he had other experiences and other endeavours than he in reality has. He is
convinced that a thought which really corresponds to a self-won conviction cannot be
altogether wrong; it will always have its significance from a certain point of view,
however limited this may be. Søren Kierkegaard developed, in the sharpest possible
form, that, as regards the great truths of life, it does not depend on what is said,
but on how it is said. It depends on the personal inwardness, on the passionate
striving and seeking, and even the most correct and truest content of faith becomes
personal untruth where this most essential condition is lacking. It is this that he
utters in his paradoxical proposition: Subjectivity is the truth\footnote{In
\emph{Udvalgte Stykker af dansk filosofisk Literatur} (1910) one will find the most
characteristic utterances of the Danish thinkers mentioned.}. It is by no means a
sign of doubt or lassitude when the strongest emphasis is thus laid on personal truth
in contrast to an objective doctrine about the worth and significance of life. It is
because among our older Danish thinkers there was a living conviction that
personality was the soil, or rather the seed, from which real wisdom of life alone can
grow forth. What for science belongs to the last objects, to whose fathoming all
general laws are to serve,---personality as a peculiar little world,---is for
life-view and faith the decisive beginning-point.

In every psychological discussion of the religious problem it will be necessary to
take account of the individual differences. I have accordingly sought, in my
``Philosophy of Religion,'' to give a survey partly of the most important oppositions
in original dispositions, partly of the differences between religious types
corresponding to these dispositions. In the first respect I find, as regards
knowledge, an opposition between natures in whom reflection and after-thought play the
chief role, and natures who hold fast as long as possible to intuition, and who,
where observation does not suffice, prefer to form images rather than give themselves
over to the movement of thought. As regards the movement of thought, there is again a
difference between natures who prefer to go into the particular, to resolve the
observations into their simplest elements,---analytic natures,---and other natures who
involuntarily strive to connect all observations and all thoughts into coherent
wholes,---synthetic natures. As regards the life of feeling and will, I find first and
foremost an opposition between natures in whom various mutually contending moods and
drives stir,---disharmonious natures, and other natures in whom the whole life of mood
and drive seeks out beyond the individual limitation in order to get an outlet and
expression for its energy,---expansive natures. There is furthermore a significant
opposition between passive and active natures,---between those whose inner life moves
in jerks or leaps, and those in whom it unfolds continuously,---and between those who
are most occupied with what stirs within themselves, and those in whom an immediate
urge to devotion prevails. When these, and surely still more typical differences,
prevail, and when they can again occur in various mutual combinations, it is clear
that both experiences and interpretations must become different. We find then also,
when we study the psychological principal forms of religious faith, very different
types, even when we hold to the same confession. I have sought to describe seven such
types. Faith can be a rest for inner and outer strife, or a free and harmonious
unfolding of the spiritual life. It can have the character of contemplation, or of
trustful venture, or of resignation. It can be an isolated individual's convulsive
holding-fast, and it can be the childlike trust in tradition of ``the charcoal-burner's
faith.'' The most important of all these oppositions is surely the one between the
disharmonious and the expansive natures,---between those in whose mind contending
forces stir, and whose urge essentially aims at finding rest, and those in whom a
great force and fullness presses forward, and whose urge therefore aims at finding
outlet and form. The life of the former will present incisive crises and
oscillations; the development of the latter will rather take place as a steadily
advancing stream. --- This opposition too has been strongly emphasized by other
psychologists of religion, and possibly all other oppositions can be led back to it.
---

As a rule, in the development of the religious consciousness, the transition from
experience to interpretation takes place very quickly. And this hangs together, as
already remarked above, with the fact that the individual finds, in the handed-down
religious representations, an interpretation laid ready in advance. Tradition, the
confession prevailing in society, does not express itself first when it is a matter of
finding form and expression, but its content weaves itself involuntarily into the
apparently immediate experiences, into visions and representations that are supposed
to arise altogether at first hand and from purely inner sources. Later the prevailing
confession works, then, when a conscious and carried-through interpretation and
presentation of what has been lived through is to be given. It has a great power in
that it is definitely formed and elaborated; it impresses by its age, and it works
perhaps even compellingly through its social organization. Therefore religious
experience consists for the most part in this, that individuals seek to rediscover and
approve the handed-down representations in their own personal life, and often they have
then forced out supposed experiences, states that in reality had no root in their
nature. It became, to use Poul Møller's expression, affectation, not personal truth.
Hence the convulsive form that religious life has often assumed, and of which the
religiosity of Pascal and Kierkegaard are such typical examples. All his genius, all
his rich life of mood, and his whole energy of will, Kierkegaard placed in the service
of the task ``to read the primal writ of the individual, human conditions of existence
--- the Old, the Familiar, and what is handed down from the Fathers --- through yet
once more.'' Life-experience thus becomes no independent and original source of faith;
the possibility and legitimacy of new interpretations are cut off in advance. The
church historian Gotfried Arnold wrote in his time (1715) a book he called
\emph{Theologia experimentalis}, and which is not merely interesting for its title. He
maintains therein very definitely that, for ecclesiastical theology, experience must
take another place than it does in science. In science experience is an independent
source and goes before faith, but in religion (confessional religion) it must come
after faith, since it is a training in what the Bible teaches, and must always ``be
judged according to God's Word.'' Arnold sees here far more clearly than many modern
theologians, who without further ado parallelize the position of experience in the
ecclesiastical and in the scientific world. There is, after all, very much in
ecclesiastical theology that no single person can have experienced, and that yet is to
be believed. It is then taken up essentially in trust in the congregation's tradition,
a kind of charcoal-burner's faith (\textit{fides implicita}). The Catholic Church is
also so liberal that it demands of laymen only a faith of this kind, a childlike trust
that the Church has the truth, even though it goes beyond what the individual can
experience or understand. It was inconsistent of Luther that he condemned a faith of
this kind, just as it was inconsistent of Kierkegaard to lay so great weight on ``the
Old, the Familiar, and what is handed down from the Fathers,'' when he at the same time
proclaimed that Subjectivity is the truth.
The free experience of life has been hampered by the preponderant role that
tradition has played in religious life. The eye has been closed to many sides of
life, of the inner life of the spirit, as well as of the outer life in its
manifoldness and extent. It is through this that the whole work done in the
direction of a scientific understanding of existence has come in advance to stand as
hostile, or in any case as indifferent, to the personal life-view. And one has been
inclined to fix a single personal life-type as the normal one, as a universal model.
There is here the contradiction, in all religion that essentially supports itself on
tradition, that it cannot explain its own beginning. When a new religion is founded,
it happens, to be sure, never on bare ground; there is always a traditional
foundation. But the tradition that constitutes the founder's spiritual atmosphere
must nevertheless undergo an essential alteration under the influence of his
personality and his life-experience. Here an insertion is made, a ``mutation''
occurs, to use a biological expression, which can contain great riddles for
historical research and psychology, but which must be acknowledged if the very
coming-to-be of the later prevailing tradition is to be understood. The difficulties
that lie herein are often veiled by the fact that one conceives the founder without
further ado, whereas his figure has first shaped itself through the successive
development of the tradition itself.

Experience must, in the domain of the life-view as in that of science, be
acknowledged as an independent source of faith, if personal truth is to come into its
own. But here there then arises a new difficulty. Science has recognized that an
assumption is not fully proved because the experiences can be understood by starting
out from it. It must also, conversely, be possible to show that from the experiences
one cannot come to another assumption, so that they can be interpreted only in this
definite way, and that no other assumption whatever is possible. It is therefore
that---as already touched upon---many researchers assume that the fundamental
principles according to which natural science has hitherto worked, and from which it
has succeeded in explaining so many phenomena in so exact a way, are not themselves
fully proved, because one cannot show the impossibility of coming to the same results
from other fundamental principles. Already Leibniz declared that a hypothesis's
success is not a proof of its truth; one must also be able to show that it is the only
possible mode of explanation. One can, after all, come to one and the same conclusion
from highly different presuppositions. If we now apply this to the personal life-view,
it is clear that it will be impossible to show that a man's experience must lead him to
an altogether definite kind of faith, whether this faith coincides with the faith of
the society in which he has grown up, or is formed anew and independently. And a
counter-test cannot here, by the nature of the case, be made. Art is long, but life is
short;---this holds also of that art of life which finds expression in a man's faith.
But it is not only time that does not permit a counter-test. In a faith or a life-view
that really is the personality's own property, comparison with other views is
excluded. It goes here as with a serious decision: the will has here concentrated
itself in such a way that all other possibilities pale; the time of comparison is past.
When the whole inwardness of the mind and the whole personal experience of life has
been deposited in our conviction, there can no longer be seriously instituted a
discussion of the question whether we might not, with our inner and outer conditions,
have come to another result. Kierkegaard comes somewhere in his diaries to touch upon
this point, in that he imagines the objection to his faith: ``Yes, but it was after all
also possible that you could have been saved in another way.'' He says with justice
that to this the believer cannot answer; if he answers the objection, it is a sign that
he in reality is not a believer. This is psychologically true; thus are we all placed,
whatever our faith may be. But from this there also follows necessarily a limitation of
all dogmatic wisdom. At no point in life is experience exhausted---and even if it were
exhausted, a conclusion drawn from it would not be absolutely compelling. It must
counsel caution. And when science itself here sees a boundary for the incontestable
validity of its assumptions and results, faith must make the same admission for its
own part. Every faith is a venture, an anticipation, and it must not settle itself so
into rest---in the individual's mind or in structures of doctrine and social
orderings---that it forgets this.

Experience and life are the broad and deep foundation on which both the scientific and
the religious concepts are built up. All our concepts are due to reflection over what
is immediately and involuntarily experienced and lived through. It is this foundation
(subject) that is gradually interpreted and determined through the utterances
(predicates) that we set forth in thoughts and in words. All our concepts are from the
first predicates. In our simplest judgments there is only a single word uttered, a
single concept formed, namely the predicate. This is the case in the pure exclamatory
judgments, for example ``Beautiful!'' ``Glorious! Great!'' The predicate is the only
thing that cannot be dispensed with in a judgment; everything else can be understood
implicitly. Often we can only point to that which the predicate holds of, without being
able to utter it. And this will especially be the case in exclamatory judgments. If now
all wisdom, the scientific as well as the religious, begins with wonder, then
exclamatory judgments will naturally belong to the first judgments. Only later can our
concepts then undergo such a determination and fixing that they can appear as subjects.
But it is of great importance to hold fast that the more we approach our original, most
virginal experiences, the more are our concepts mere predicate-concepts. This is their
first stage, to which every thorough discussion must go back. And especially the
religious concepts constantly need a revision in this direction. The very concept God
is thus a predicate-concept, has arisen as a thought-expression for experiences and
lived events. It becomes clear as soon as one asks what makes God God or conditions
God's divinity. It then shows itself that in the concept God, at every time and for
every kind of religion, are comprised the most valuable properties that men found made
known in existence and in the march of destiny, such as they knew and understood it. It
is the common trait beneath the infinite variation of the representations of God. And
it is to that the religious discussion will always at last come back, however difficult
it may be to discover a concept's real hour of birth. The question for the individual
will be whether he, according to his experience, can repeat the process by which the
handed-down religious representations originally came to be. Possibly these
representations must be considerably altered if his experience is to approve them.
Possibly he will come to form wholly other representations, if he is to be true to his
real experience. It will perhaps be impossible for him to form altogether definite
representations, or he will become clear that only images and parables can here help
him.

We stand here at a chief point of view for the religious problem in our days. On many
points the strife will, as hitherto, be carried on from the side of natural science, of
historical criticism, and of critical philosophy. But science and philosophy cannot
bring forth faith and are not enough for a personal life-view. Faith and life-view grow
involuntarily out of life, and thinking plays more a role as a limiting and testing
power than as the positively producing cause. Science can only demand that, just as
little as it claims to be able to fathom existence, just as little must religion suppose
that it could ever provide universally necessary and exhaustive expressions and forms
for the experiences of personal life. There must constantly be a distinction between
experience and interpretation. It is this distinction that Danish thinkers have wished
to inculcate with their demand for personal truth, and whose importance the newest
psychological research inculcates. It is now no longer merely a matter of interpreting
old books or confessions, but of interpreting life itself, the life that is richer and
deeper than any tradition can express. If along this path new wine is harvested, it
will show itself whether it can be contained in the old wineskins.

The recognition of religious freedom has brought a step forward in the direction of
opening the possibility of a free personal life. Religious freedom rests on the
presupposition that also outside the ecclesiastically closed religious communities a
worthy and noble human life can be led. First it was wrung from the ecclesiastical
power-holders; but it is becoming a positive principle, an expression of the trust in
the independent significance of human experience and human development. The more the
Church gives up working with outer power and as a state institution---that is, in the
course of its progressive separation from the State---the more room and the more
necessity there is for independent life-experience. The Church itself will come to work
for its ideals under other conditions than before; it now steps, for the first time
after its heroic age, down onto the arena and is placed on an equal footing with other
human endeavours for what can make life great and beautiful. And it will then be
obliged to reflect more than before, and notably with greater clarity, upon the
foundation of experience on which it can build. As long as it was its own traditions
that, in more or less marked form, determined the whole spiritual atmosphere in which
all individuals grew up, so long its work was not so difficult. It consisted for a great
part in getting forgotten childhood melodies and homely memories to become living again.
Now it will be obliged to build its edifice up from the ground. And as regards the
Church's opponents, the conditions thus altered will gradually bring the good, that
criticism of the handed-down will not take up so much of the independently striving
ones' time and force. In all too many there were awakened and nourished, under the
overlordship of the religious traditions, a fear and a reluctance toward every positive
life-view, and one had grown so accustomed to a critical attitude that one forgot all
criticism's merely preparatory and provisional significance. Now indifference or lack of
trust in life's positive unfolding in new forms, now downright blaséness, became the
consequence. There will perhaps go a long time before the unhappy consequences of
dogmatism will be overcome. But it is, after all, an old saying that he who believes is
in no hurry.

Perhaps the strife will become eternal between those whose thoughts about life can be
formed in a confession they have in common with others on the ground of tradition, and
those whose inner life---whether it now be an advantage or a defect---cannot be laid to
rights in this way. A spirited Danish clergyman once said: ``A seeking congregation is
an absurdity.'' It is true that a religious community must join together about a definite
conception of life, and that the constant seeking must become the individuals' affair.
There are also many historians of religion who so highly emphasize that religion is a
social phenomenon, that they take no interest in the more individually marked forms of
faith, which yet are both beginning-points and high-points in the history of religion.
But from the point of view we have here to do with, the question becomes whether life
always stirs more strongly in those who have all the dogmatic rubrics filled in, than in
those who, despite serious seeking, do not land in any concluding confession. If the
inmost, the most personal, and the most valuable that can be lived and experienced could
not be expressed in any dogma, as little as in any scientific proposition, was it then an
absurdity to be a seeker and to continue to be one?

This question is the most important that can be posed in our days. Its answer forms the
decisive dividing-line. It has its significance in laying upon our minds that in life as
in science it is in the end a matter of building on experiences one has oneself made, and
whose interpretation one does not take from the hand of others. The individual develops,
both in life and in science, under the influence of the traditions of the race, both the
venerable and the questionable. They are for him now a help, now a barrier. But whether
in the particular case they are worthy of honour or not, and whether they are help or
barrier, that can for the individual in the end only be decided by his own personal
experience.

The strife in the world of the spirit will be carried on under better conditions than
hitherto, if this way of viewing things presses more and more through. And only if that
happens can he who cares both about the freedom of thought and about the richness and
inwardness of personal life look toward the future with confidence.

\begin{center}---\end{center}

\end{document}
